\section{What has changed, what remains, and the next gate}
\label{sec:tpc48-boundaries}

The present step closes the geometric all-band question raised in
TPC--47.  It does not close the arithmetic fixed-shift question.  We
record the distinction as a theorem ledger before describing a
possible continuation.

\subsection{Advances under the inherited exact interface}

The following statements are proved here without a new arithmetic
hypothesis beyond the exact TPC--36/45/47 representation and support
interface.
\begin{enumerate}[label=\textup{(\roman*)},leftmargin=2.4em]
 \item The orbit-frequency quotient admits an exact orthogonal tiling
       by \(L\) full residue combs, with exact compressed trace \(N/L\)
       on every scalar physical window.
 \item Each determinant modulation has an exact input-output channel
       resolution.  The support consists of explicit resonance tubes,
       and the critical branch half-width is \(X^{-1+o(1)}\).
 \item The optimal norm of the channel lift is exactly the maximal
       pointwise tube multiplicity.  Erasing the channel label restores
       the unrestricted synthesis norm squared \(K\).
 \item The common orbit profile in every inherited projective layer
       yields the Hilbert-valued factor
       \(D_J(t;K)=\min\{K,1+(Jt)^{-1}\}\), sharp up to constants in its
       declared class.
 \item Complete signed reassembly, coefficient by coefficient, reduces
       the literal frozen family to a representation independent
       pointwise
       projective \(m\)-square envelope.  At fixed \(m\), residual
       labels give an exact outer \(r\)-square sum, while the prime sum
       inside each residual coordinate remains coherent.
 \item Relative to any fixed global projective representation with a
       bounded envelope, the intermediate annulus has only logarithmic
       integrated resonance cost.
\end{enumerate}

These are forward results rather than a reformulation of the
single-band theorem: all orbit bands are now controlled at once, the
fine determinant scale is restored, and the exact obstruction left by
their reassembly is identified.

\subsection{The remaining arithmetic square-function gate}

By \eqref{eq:tpc48-residual-square}, the unresolved layerwise object is
\begin{equation}
 \sum_m\sum_r
 \left\|
 \sum_p A_\omega(p,r,m)\e(-pr\alpha)
 v_{\omega,p,r,m}
 \right\|^2.
 \label{eq:tpc48-next-arithmetic-object}
\end{equation}
Three features prevent the present large sieve from estimating this
quantity at the atomic scale:
\begin{enumerate}[label=\textup{(\alph*)},leftmargin=2.4em]
 \item the \(p\)-sum is coherent inside a fixed residual coordinate;
 \item the projective integral must be reassembled with its literal
       signs and total variation, not by selecting a favorable layer;
 \item subsequent residual grouping can cost
       \(\rho_{\rm res}^{1/2}\), for which no \(X^{o(1)}\) bound has
       been inherited at this interface.
\end{enumerate}
For the dense consecutive-slope sharpness model at the coarse edge,
the common-profile collision factor exceeds the TPC allowance by the
exponent \(\nu=133/400\).  Any continuation
that aims at atomic closure must therefore use arithmetic information
inside \eqref{eq:tpc48-next-arithmetic-object}; orbit geometry alone is
now known to be insufficient.

\subsection{Four precise doors}

The next paper may attack any one of the following independently.
\begin{description}[leftmargin=2.8em,style=nextline]
 \item[Door A: coefficient-specific prime square.] Prove cancellation
   in the coherent prime sum of
   \eqref{eq:tpc48-next-arithmetic-object}, averaged over \(m,r\),
   beyond the universal common-profile large sieve.
 \item[Door B: projective-to-atomic comparison.] Bound one fixed exact
   representation envelope \(\mathfrak S_{\mathscr R}\) by the
   inherited physical atomic diagonal with a quantified loss, or
   construct an actual-mask witness showing that such a bound is
   impossible at the desired exponent.
 \item[Door C: residual grouping with cofactor structure.] Replace the
   abstract \(\rho_{\rm res}\) loss by a coefficient-specific
   square-function estimate using the literal map \(s=d(m)r\).  The
   balanced cofactor fibers from earlier descent cannot be imported
   here without proving that the two interfaces coincide.
 \item[Door D: localization and alias completion.] Combine the
   annular theorem with an ultra-low analysis, an anti-alias or
   maximal-in-shift theorem, and the literal \(h_0\) Fourier inversion.
   This door is downstream of, rather than a substitute for, the
   atomic comparison.
\end{description}

Each door has a meaningful negative outcome: an explicit actual-mask
counterexample or a sharp lower bound for the necessary loss would
locate the next obstruction just as the unrestricted synthesis theorem
does here.

\subsection{Prohibited conclusions}

For clarity, none of the following is asserted.
\begin{enumerate}[label=\textup{(\roman*)},leftmargin=2.4em]
 \item No estimate of the complete physical atom by its atomic
       diagonal is proved.
 \item No bound \(\rho_{\rm res}=X^{o(1)}\), no new cofactor-range
       closure, and no cancellation of the internal prime sum is
       proved.
 \item No sparse-alias unfolding, new-shift physical provenance, or
       uniform local-shift theorem is proved.  Only \(h=h_0\) is the
       inherited physical residual synthesis; all other shifts are
       frozen algebraic extensions.
 \item The residual synthesis remains the low-sign-averaged Gram
       route and is not identified with the native affine physical
       face.
 \item No parity-barrier breach, Hardy--Littlewood asymptotic, bounded
       prime-gap theorem, or twin-prime conclusion follows.
\end{enumerate}

The exact conclusion is narrower and reusable: after full
orbit-frequency resolution, the next non-geometric obstruction is the
actual projective \(m,r\)-square of coherent prime sums.  That is the
next gate, not a disguised solution of it.
